\section{Preliminaries}\label{sec:pre}
This section provides the precise definitions of density and measurability.
We denote by \(\reg_A\) the family of all regular languages over an alphabet \(A\). For a word \(w \in A^*\) and a letter \(a \in A\), we write \(|w|_a\) for the number of occurrences of \(a\) in \(w\). The complement of \(L\) over \(A\) is denoted by \(\comp{L} = A^* \setminus L\).
We assume that the reader has a standard knowledge of algebraic language theory (\cf \cite{Lawson,DAM,MPRI}).

For a finite monoid \(M\) and an element \(x \in M\), the unique idempotent power of \(x\) is denoted \(x^\omega\).
A moniod \(M\) is called \emph{aperiodic} if \(x^{\omega+1} = x^\omega\) holds for any \(x \in M\).
Recall that Green's \(\CH\)-relation on a monoid \(M\) is defined as \(x \CH y \Leftrightarrow xM = yM \land Mx = My\), and
\(\CJ\)-relation is defined as \(x \CJ y \Leftrightarrow MxM = MyM\).
A monoid \(M\) is said to be \({\cal G}\)-trivial if every \({\cal G}\)-class is a singleton where \({\cal G} \in \{\CH, \CJ\}\).
It is well-known that a finite moniod \(M\) is aperiodic if and only if \(M\) is \(\CH\)-trivial.
The main targets of this paper are the following three subclasses of regular languages. The \emph{star-free languages}, the \emph{generalised definite languages} \(\gd\), and the \emph{group languages} \(\group\):
\begin{align*}
 \sf_A &\defeq \set{L \subseteq A^*}{L \text{ is recognised by a finite \(\CH\)-trivial monoid}},\\
 \gd_A &\defeq \bool\set{wA^*, A^*w}{w \in A^*},\\
 \group_A &\defeq \set{L \subseteq A^*}{L \text{ is recognised by a finite group}}.
\end{align*}
Also, we denote by \(\gcom_A\) the class of all languages recognised by finite commutative groups.
A monoid \(M\) divides a monoid \(N\) if \(M\) is a monoid homomorphic image of a submonoid of \(N\).
It is well-known that a monoid \(M\) recognises a language \(L\) if and only if the syntactic monoid \(M_L\) of \(L\) divides \(M\), hence
\(L\) is in \(\sf_A\) (\resp \(\group_A\)) if and only if \(M_L\) is a finite \(\CH\)-trivial monoid (\resp a finite group).
It is easy to see that a language \(L\) is generalised definite if and only if \(L = E \cup \bigcup_{(u, v) \in F} uA^*v\) for some finite sets \(E \subseteq A^*\) and \(F \subseteq A^* \times A^*\).
% \(L = E \cup \bigcup_{i \in I} F_i A^* G_i\) for some finite sets \(E\) and \(F_i, G_i \subseteq A^*\) for all \(i \in I\), where \(I\) is a finite index set.

\subsection{Density and measurability of formal language}
For a set $X$, we denote by $\card{X}$ the cardinality of $X$.
We denote by $\nat$ the set of natural numbers including $0$.

\begin{definition}[\cf \cite{codes}]\upshape
The \emph{density} \(\cd(L)\) of \(L \subseteq A^*\) is defined as
\[
 \cd(L) \defeq \lim_{n \rightarrow \infty} \frac{1}{n} \sum_{k = 0}^{n-1} \frac{\card{L \cap A^k}}{\card{A^k}}
\]
 if it exists, otherwise we write \(\cd(L) = \bot\).
 A language \(L\) is called \emph{null} if \(\cd(L) = 0\),
 and dually, \(L\) is called \emph{co-null} if \(\cd(L) = 1\).
\end{definition}

\begin{example}\label{ex:density}
It is known that every regular language has a rational density (\cf \cite{codes}) and it is computable.
For each word \(w\), the language \(A^* w A^*\),
 the set of all words that contain \(w\) as a factor,
 is of density one (\ie, co-null). This fact follows from the so-called {\em infinite monkey theorem}:
take any word \(w\). A random word of length \(n\) contains \(w\) as a factor with probability tending to one as \(n\) tends to infinity.
\end{example}

We list some basic properties of the density as follows.
\begin{lemma}[\cite{codes}]\label{lem:density}
 Let \(K, L \subseteq A^*\) with \(\cd(K) = \alpha, \cd(L) =
 \beta\). Then we have:
\begin{enumerate}
\item \(\cd(L \setminus K) = \beta - \alpha\) if \(K \subseteq L\). 
% \item \(\cd(\comp{K}) = 1 - \alpha\).
% \(\cd(K \cup L) \leq \alpha + \beta\) if \(\cd(K \cup L) \neq \bot\). 
\item \(\cd(K \cup L) = \alpha + \beta\) if \(K \cap L = \varnothing\). 
\item \(\cd(uLv) = \cd(L) \cdot \card{A}^{-|uv|}\) for each \(u, v \in A^*\).
\end{enumerate}
\end{lemma}
\begin{lemma}[{\cite[Example 13.4.9]{codes}}]\label{lem:groupdensity}
 Let \(\eta\colon A^* \rightarrow G\) be a surjective morphism onto a finite group. For any \(g \in G\), we have \(\cd(\eta^{-1}(g)) = \card{G}^{-1}\).
\end{lemma}

\begin{remark}\label{rem:kernel}
Let \(L\) be a regular language, \(M\) be a finite monoid and 
\(\eta: A^* \rightarrow M\) be a surjective morphism such that \(L = \eta^{-1}(\eta(L))\).
The \emph{kernel} \(K\) of \(M\) is a unique minimal ideal of \(M\) (it is always exists since \(M\) is finite),\ie, the minimal \(\CJ\)-class of \(M\).
For every \(m\) not in the kernel \(K\) of \(M\),
\(\eta^{-1}(m) \cap A^* w A^* = \emptyset\) holds for any \(w \in \eta^{-1}(K)\) hence \(\d(\eta^{-1}(m)) = 0\) by the infinite monkey theorem.
This implies that \(\d(\eta^{-1}(K)) = 1\).
Moreover, \(\d(\eta^{-1}(K)) = 1\) implies that there exist at least one element \(n \in K\) such that \(\d(\eta^{-1}(n)) > 0\) since \(K\) is finite.
Because \(K\) is the kernel of \(M\), for any \(m \in K\) there exists a pair of 
element \(x, y \in M\) such that \(xny = m\), hence \(\eta^{-1}(m)\) contains \(u \eta^{-1}(n) v\) where \(u \in \eta^{-1}(x)\) and \(v \in \eta^{-1}(y)\) so that
\(\d(\eta^{-1}(m)) \geq \d(\eta^{-1}(n)) \cdot \card{A}^{-(|u|+|v|)}\) by Lemma~\ref{lem:density}, \ie, \(\d(\eta^{-1}(m)) > 0\) holds for each \(m \in K\).
\end{remark}
For more detailed properties of the density function \(\cd\), see Chapter~13 of \cite{codes}.

The notion of ``measurability'' on formal languages is defined by a standard measure theoretic way as follows.
\begin{definition}[\cite{S2021}]\upshape
Let \(\CC_A\) be a family of languages over \(A\).
For a language \(L \subseteq A^*\), we define its
 \emph{\(\CC_A\)-inner-density} \(\plm{\CC_A}(L)\)
 and \emph{\(\CC_A\)-outer-density} \(\pum{\CC_A}(L)\)
 over \(A\) as
 \[
 \begin{alignedat}{2}
  \plm{\CC_A}(L) &\defeq& \, \sup &\set{\cd(K)}{K \subseteq L, K \in \CC_A, \cd(K) \neq \bot} \text{ and }\\
  \pum{\CC_A}(L) &\defeq& \inf &\set{\cd(K)}{L \subseteq K, K \in \CC_A, \cd(K) \neq \bot} \text{, respectively.}
 \end{alignedat}
 \]
A language \(L\) is said to be
 \emph{\(\CC_A\)-measurable} if
 \(\plm{\CC_A}(L) = \pum{\CC_A}(L)\) holds.
 We say that an infinite sequence \((L_n)_n\) of languages over \(A\)
 {\em converges to \(L\) from inner ({\em \resp}from outer)} if
 \(L_n \subseteq L\) (\resp \(L_n \supseteq L\)) for each \(n\) and \(\lim_{n \rightarrow \infty}\cd(L_n) = \cd(L)\).
A language \(L\) is \(\CC_A\)-measurable if and only if
there exist two convergent sequences of languages in \(\CC_A\) to \(L\) from inner and outer, respectively.
\end{definition}

The following is an example of \(\gcom\)-measurable non-regular languages.
For more detailed examples of \(\CC\)-measurable/immeasurable languages for a subclass \(\CC\) of regular languages, see \cite{S2021,carathe,gcommeasure}.

\begin{example}[\cite{S2021}]\label{ex:dyck2}
 The semi-Dyck language \(\dyck = \{\varepsilon, ab, aabb, abab, aaabbb, \ldots\}\) over \(A = \{a,b\}\) is \(\gcom\)-measurable.
 For each \(k \geq 2\), the language \(L_k = \{ w \in A^* \mid |w|_a = |w|_b \mod k \}\) is in \(\gcom\), because 
 \(\eta^{-1}(0) = L_k\) holds for the morphism \(\eta\colon A^* \rightarrow \modint{k}\) where \(\eta(a) = 1\) and \(\eta(b) = k-1\).
 Obviously, \(\dyck \subseteq L_k\) holds and it follows from Lemma~\ref{lem:groupdensity} that
 \(\cd(L_k) = 1/k\) holds. Hence \(\cd(L_k)\) tends to zero if \(k\) tends to infinity.
\end{example}

\subsection{Local varieties and profinite equations}
For a family \(\CC_A\) of languages over \(A\), we denote by
\(\ext(\CC_A)\) (\resp \(\rext(\CC_A)\)) the class of all \(\CC_A\)-measurable languages (\resp \(\CC_A\)-measurable \emph{regular} languages) over \(A\).
A family of regular languages over \(A\) is called a \emph{local variety} (\cf \cite{Adamek}) over \(A\) if it is closed under 
 Boolean operations and left-and-right quotients.
All classes \(\gd_A, \sf_A, \gcom_A\), and \(\group_A\) are local varieties.
The notion of \(\CC_A\)-measurability is well-behaved on Boolean operations and quotients as the following lemma says.
% A family of regular languages over \(A\)
%  is called a {\em local variety}~\cite{Adamek} over \(A\) if it is closed under 
%  Boolean operations and left-and-right quotients.
\begin{lemma}[\cite{carathe}]\label{lem:closure}
The operator \(\ext\) is a closure,
 \ie, it satisfies the following three properties for each \(\CC_A \subseteq \CD_A \subseteq 2^{A^*}\):
 (extensive)  \(\CC_A \subseteq \ext(\CC_A)\),
 (monotone) \(\ext(\CC_A) \subseteq \ext(\CD_A)\) and
 (idempotent) \(\ext(\ext(\CC_A)) = \ext(\CC_A)\).
Moreover, 
if \(\CC_A\) is a local variety, then \(\ext(\CC_A)\) is also a local variety.
\end{lemma}

Let \(\fin_A\) be the class of all finite and co-finite languages over \(A\).
Hereafter, we only consider an alphabet \(A\) such that \(\card{A} \geq 2\) (because \(\sf_A = \gd_A = \fin_A\) and \(\group_A = \gcom_A\) holds if \(\card{A} = 1\)) and sometimes omit the subscript \(A\) for denoting these classes of languages.

The notion of local varieties and local pseudovarieties can be captured by an appropriate ``equation'' notion called \emph{profinite equations} (\cf \cite{profinite}).
For a pair of words \(u\) and \(v\), we can define a dinstance between them by
\(d(u, v) = 2^{-r(u,v)}\)
where \(r(u,v) = \min\{ |\CA| \mid \CA \text{ is a deterministic automaton accepting one of } \{u, v\} \text{ but not the other } \}\).
Intuitively, \(r(u,v)\) represents the minimum number of states for separating \(u\) and \(v\) by using deterministic automata.
Then this function \(d\) is an ultrametric on \(A^*\), and the completion of the metric space \((A^*, d)\), denoted by \(\widehat{A^*}\), is called the set of profinite words over \(A\).
An element of \(\widehat{A^*}\) is a Cauchy sequence of (usual finite) words over \(A\) with respect to \(d\).
For example, the profinite word \(w^\omega\), associated with every word \(w\), can be defined as a limit of the sequence of words \((w^{n!})_{n \geq 1}\).
For any finite monoid \(M\) and morphism \(\eta: A^* \rightarrow M\), it is easy to see that \(\eta(w^{n!})\) is the idempotent power \(\eta(w)^{\omega}\) of \(\eta(w)\).
Hence, for a finite monoid \(M\), the equation \(w^{\omega+1} = w^{\omega}\) can be interpreted as \(x^{\omega+1} = x^\omega\) holds for any \(x \in M\), \ie, 
\(M\) is aperiodic (= \(\CH\)-trivial).
It is known that any local variety of languages can be captured by using profinite equations and vice versa.
We use the following theorem to prove Theorem~\ref{thm:group-robust} in Section~\ref{sec:group}.
%For the definition and basic properties of the free profinite monoid \(\widehat{A^*}\), see \cite{profinite,MPRI} for example.
\begin{theorem}[{\cite[Proposition 7.4]{Duality2008}}]\label{thm:duality}
  A set of regular languages over \(A\) is a Boolean quotienting algebra (\ie, local variety)
  if and only if it can be defined by a set of semigroup equations of the form \(u = v\), where \(u, v \in \widehat{A^*}\).
\end{theorem}
A regular language \(L \subseteq A^*\) \emph{satisfies} an equation \(u = v\), where \(u, v \in \widehat{A^*}\),
if the syntactic morphism \(\eta\colon A^* \to M_L\) of \(L\) satisfies \(\hat{\eta}(u) = \hat{\eta}(v)\),
where \(\hat{\eta}\colon \widehat{A^*} \to M_L\) is the unique uniformly continuous extension of \(\eta\)
(\cf \cite[Chapter X, Proposition 2.10]{MPRI}).

\begin{lemma}\label{lem:group-profinite-identity}
  For any equation \(u = v\) (where \(u, v \in \widehat{A^*}\)),
  there exists \(w \in \widehat{A^*}\) such that the equation \(w = \varepsilon\) is equivalent to \(u = v\) on any \(A\)-generated finite group.
\end{lemma}
\begin{proof}
  Let \((v_n)_{n \in \nat}\) be a Cauchy sequence in \(A^*\) converging to \(v\) in \(\widehat{A^*}\).
  The profinite word \(v^{\omega - 1} \in \widehat{A^*}\) is defined as \(v^{\omega - 1} = \lim_{n \to \infty} v_n^{n! - 1}\).
  Then, for every \(A\)-generated finite group \(\eta\colon A^* \to G\), we have \(\hat{\eta}(v^{\omega - 1}) = \hat{\eta}(v)^{-1}\) in \(G\).
  Thus \(w = u v^{\omega - 1}\) satisfies the condition of the lemma.
\end{proof}

\subsection{Local varieties of finite nilpotent/solvable groups}
This subsection provides background material for Section~\ref{sec:group}.
The reader may skip this subsection if interested only in the results up to Section~\ref{sec:ind}.

\begin{definition}
  For a positive integer \(n\), define
  \begin{align*}
    \gnil_A^{\leq n} &\defeq \set{L \subseteq A^*}{\text{\(L\) is recognised by a finite nilpotent group of class \(\leq n\)}},\\
    \gsol_A^{\leq n} &\defeq \set{L \subseteq A^*}{\text{\(L\) is recognised by a finite solvable group of derived length \(\leq n\)}}.
  \end{align*}
  Note that \(\gnil_A^{\leq 1} = \gsol_A^{\leq 1} = \gcom_A\).
\end{definition}
Recall that a group \(G\) is called a nilpotent group of class \(\leq n\) if and only if for any \(g_0, g_1, g_2, \dotsc, g_n \in G\),
the equation \([[\dotsm[[g_0, g_1], g_2]\dotsm], g_n] = 1_G\) holds (\eg, \cite[14.64]{hanna-neumann-variety}),
where \([g, h] \defeq g h g^{-1} h^{-1}\) is the commutator of \(g\) and \(h\).
Hence, for example, \(\gnil_A^{\leq 2}\) is a local variety defined by the set of equations of the form
\((u v u^{\omega - 1} v^{\omega - 1}) w (u v u^{\omega - 1} v^{\omega - 1})^{\omega - 1} w^{\omega - 1} = \varepsilon\) for \(u, v, w \in A^*\).
Similarly, each \(\gnil_A^{\leq n}\) is also a local variety.
It is also known that each \(\gsol_A^{\leq n}\) is defined by a set of equations (\eg, \cite[14.65]{hanna-neumann-variety}),
therefore \(\gsol_A^{\leq n}\) is also a local variety.

%For each \(n \geq 1\), it is known that there exists a finite nilpotent group of class \(n\) that is generated by at most \(n\) elements.
%For example, if \(p\) is a positive integer,
%then the group \(\UpTriMat(n, \integer / p \integer)\) of \(n \times n\) unitriangular matrices over the ring \(\integer / p \integer\)
%is a finite nilpotent group of class \(n - 1\), generated by \(n\) elements~\cite[Exercise 5.44.(ii)]{rotman-GTM148}.
%Therefore \(\gnil_A^{\leq n} \subsetneq \gnil_A^{\leq n + 1}\) for each \(n < \card{A}\).

%For each \(n \geq 1\), it is also known that there exists a finite solvable group of derived length \(n\)
%that is generated by at most \(n\) elements~\cite[9.22 Lemma]{rose-group}.
%For example, the \(n\)-fold iterated wreath product \((\integer / 2 \integer) \wr \dotsb \wr (\integer / 2 \integer)\)
%is a finite solvable group of derived length \(n\), generated by \(n\) elements.
%Therefore \(\gsol_A^{\leq n} \subsetneq \gsol_A^{\leq n + 1}\) for each \(n < \card{A}\).

%\begin{lemma}\label{lem:two-gen}
%  Any finite group \(G\) is generated by \(n\) elements as a group if and only if it is generated by \(n\) elements as a monoid.
%  In particular, if \(G\) is generated by two elements, there exists a surjective morphism \(\eta\colon A^* \to G\).
%\end{lemma}
%\begin{proof}
%  Since \(G\) is finite, the inverse of each \(g \in G\) is expressed as \(g^{\card{G} - 1}\).
%\end{proof}

The rest of this subsection is devoted to the proof of Proposition~\ref{prop:strict}.
\begin{proposition}\label{prop:strict}
  Suppose \(\card{A} \geq 2\).
  For every \(n \geq 1\), we have \(\gnil_A^{\leq n} \subsetneq \gnil_A^{\leq n + 1}\) and \(\gsol_A^{\leq n} \subsetneq \gsol_A^{\leq n + 1}\).
\end{proposition}
To prove Proposition~\ref{prop:strict}, we need Lemmata~\ref{lem:two-gen-nilpotent} and \ref{lem:two-gen-solvable} below.

\begin{lemma}\label{lem:two-gen-nilpotent}
  For \(n \geq 1\), there is a two-generated finite nilpotent group of class \(n\).
\end{lemma}
\begin{proof}
  For each \(n \geq 1\), the dihedral group \(D_{2^{n + 1}}\) of order \(2^{n + 1}\)
  is a two-generated finite nilpotent group of class \(n\) (see \eg, \cite[Chapter 4, \S 4]{suzuki-2-ja}).
\end{proof}

\begin{lemma}\label{lem:two-gen-solvable}
  For \(n \geq 1\), there is a two-generated finite solvable group of derived length \(n\).
\end{lemma}
To prove Lemma~\ref{lem:two-gen-solvable}, we use the following embedding theorem due to Neumann-Neumann,
and Lemmata~\ref{lem:solvable-length-n} and \ref{lem:solvable-reduce} below.

\begin{theorem}[{\cite[Theorem 5.6]{neumann-neumann-1959}}]\label{thm:embed-solvable}
  Every finite solvable group \(G\) of derived length \(n\) can be embedded in
  some two-generated finite solvable group \(H\) of derived length at most \(n + 2\).
\end{theorem}

%\begin{lemma}[{\cite[Lemma 2.1]{AM2011}}]\label{lem:length-wreath}
%  Let \(G\) be a solvable group of derived length \(n\) and \(k \geq 2\) be an integer.
%  Then the wreath product \(G \wr C_k\) is a solvable group of derived length \(n + 1\),
%  where \(C_k\) is the finite cyclic group of order \(k\).
%\end{lemma}

\begin{lemma}\label{lem:solvable-length-n}
  For each \(n \geq 1\), there exists a finite solvable group of derived length \(n\).
\end{lemma}
\begin{proof}
  %By Lemma~\ref{lem:length-wreath},
  %the \(n\)-fold iterated wreath product \((\integer / 2 \integer) \wr \dotsb \wr (\integer / 2 \integer)\)
  %is a finite solvable group of derived length \(n\).
  For example, the \(n\)-fold iterated wreath product \((\integer / 2 \integer) \wr \dotsb \wr (\integer / 2 \integer)\)
  is a finite solvable group of derived length \(n\)~\cite[9.22 Lemma]{rose-group}.
\end{proof}

\begin{lemma}\label{lem:solvable-reduce}
  If \(G\) is a solvable group of derived length \(n + 1\), then some factor group of \(G\) is solvable of derived length \(n\).
\end{lemma}
\begin{proof}
  Let \(G = G^{(0)} \rhd G^{(1)} \rhd \dotsb \rhd G^{(n)} \rhd G^{(n + 1)} = \{1_G\}\) be the derived series of \(G\),
  where \(G^{(i + 1)} \defeq [G^{(i)}, G^{(i)}]\).
  Since each \(G^{(i + 1)}\) is a characteristic subgroup of \(G^{(i)}\), the group \(G^{(n)}\) is also a characteristic subgroup of \(G\),
  hence in particular a normal subgroup of each \(G^{(i)}\).
  Let \(\bar{G} \defeq G / G^{(n)}\).
  Then \(G^{(0)} / G^{(n)} \rhd G^{(1)} / G^{(n)} \rhd \dotsb \rhd G^{(n)} / G^{(n)} = \{1\}\) is the derived series of \(\bar{G}\) of length \(n\),
  since we have an isomorphism \([G^{(i)} / G^{(n)}, G^{(i)} / G^{(n)}] \cong [G^{(i)}, G^{(i)}] / G^{(n)} = G^{(i + 1)} / G^{(n)}\).
\end{proof}

\begin{proof}[Proof of Lemma~\ref{lem:two-gen-solvable}]
  By Lemma~\ref{lem:solvable-length-n}, there exists a finite solvable group \(G\) of derived length \(n\).
  By Theorem~\ref{thm:embed-solvable},
  \(G\) can be embedded in some two-generated finite group \(H\) of derived length at most \(n + 2\).
  %By Lemma~\ref{lem:length-wreath}, the derived length of \(G \wr C_k\) is \(n + 1\).
  In particular, the derived length of \(H\) is one of \(n\), \(n + 1\), or \(n + 2\).
  Thus, by Lemma~\ref{lem:solvable-reduce},
  we obtain a two-generated finite solvable group of derived length \(n\) as a factor group of \(H\).
\end{proof}

Lemma~\ref{lem:WP-synt} below is the last piece of the proof of Proposition~\ref{prop:strict}.
%Let \(G\) be a group and \(\eta\colon A^* \to G\) be a surjective morphism.
%The \emph{word problem} \(\WP_\eta(G)\) of \(G\) (with respect to \(\eta\))
%is the language over \(A\) defined as \(\WP_\eta(G) \defeq \eta^{-1}(1_G)\).
\begin{lemma}[{\cite[3.4.9 Proposition]{Holt2017}}]\label{lem:WP-synt}
  Let \(G\) be a group and \(\eta\colon A^* \to G\) be a surjective morphism.
  Then the syntactic monoid of the language \(\eta^{-1}(1_G) \subseteq A^*\) is isomorphic to \(G\).
  %Then the syntactic monoid of \(\WP_\eta(G)\) is isomorphic to \(G\).
  %In particular, every \(A\)-generated group is the syntactic monoid of some language over \(A\).
\end{lemma}
The language \(\eta^{-1}(1_G)\) is known as the \emph{word problem} for the group \(G\) (over \(A\)).

\begin{proof}[Proof of Proposition~\ref{prop:strict}]
  By Lemma~\ref{lem:two-gen-nilpotent}, there exists a two-generated finite nilpotent group \(G\) of class \(n + 1\).
  Since \(G\) is finite and generated by two elements as a group, it is generated by the two elements as a monoid.
  Hence there exists a surjective morphism \(\eta\colon A^* \to G\).
  Since \(\eta^{-1}(1_G)\) is clearly recognised by \(G\), we have \(\eta^{-1}(1_G) \in \gnil_A^{\leq n + 1}\).
  Contrary, suppose \(\eta^{-1}(1_G) \in \gnil_A^{\leq n}\).
  Then \(\eta^{-1}(1_G)\) is recognised by some finite nilpotent group \(H\) of class \(n\),
  hence the syntactic monoid of \(\eta^{-1}(1_G)\), which is isomorphic to \(G\) by Lemma~\ref{lem:WP-synt}, divides \(H\).
  In general, if \(K\) is a nilpotent group of class \(c\),
  then every subgroup of \(K\) and every factor group of \(K\) is nilpotent of class at most \(c\).
  Hence \(G\) is nilpotent of class at most \(n\), contradiction.
  (Note that every submonoid of a group is a subgroup.)
  Thus \(\gnil_A^{\leq n} \subsetneq \gnil_A^{\leq n + 1}\).

  By Lemma~\ref{lem:two-gen-solvable}, there exists a two-generated finite solvable group \(G\) of derived length \(n + 1\).
  Then one can similarly show \(\eta^{-1}(1_G) \in \gsol_A^{\leq n + 1} \setminus \gsol_A^{\leq n}\) for a surjective morphism \(\eta\colon A^* \to G\).
\end{proof}

